\section{Glossary of terms and hypotheses}\label{app:definitions}

This appendix collects terms used across several sections. It introduces no new mathematics; the
controlling definitions are the ones cited.

\subsection{Terms}\label{subsec:definitions-list}

\begin{description}[style=multiline,leftmargin=!,labelwidth=0.27\textwidth,labelsep=0.8em,itemsep=3pt]
\item[Theory package.] A carrier, lens, packaging map and audit $(Z,f,E,\mathcal A)$
(Convention~\ref{conv:package}), used here as an organizing instrument.
\item[Presentation.] A finite pointed acyclic graph $(V,\to,v)$ with $V\subseteq\mathbb N$ finite,
read as a membership presentation (Definition~\ref{def:presentation}). The set $Z$ of all
presentations is countably infinite.
\item[Membership-profile lens.] The map $f\colon Z\to\HF$ sending a presentation to the decoration
of its point (Definition~\ref{def:lens}).
\item[Packaging map.] $E(P)=C(f(P))$, which forgets the shape of a presentation and re-presents its
profile canonically (Definition~\ref{def:E}).
\item[Rank audit.] $\rho(P)=\rank(f(P))$, strictly decreasing along membership
(Definition~\ref{def:rank}, Lemma~\ref{lem:L3}).
\item[Finite packaging operator.] The operator $F$ on subfamilies of $\HF$ generated by the empty
set, pairing, union, subsets (gates) and finite images (transport) (Definition~\ref{def:F}).
\item[Finitely generated rule.] A completion rule adding only finite sets whose members come from
the current family and the transitive closures of its members (Definition~\ref{def:fg}).
\item[Rank-erasure.] Passing from the two-sorted staged structure to the one-sorted membership
structure (Definition~\ref{def:staged}).
\item[Coded predicates $\mathcal B_M$.] $\powset(\omega)\cap M$, a countable Boolean algebra of
predicates on $\omega$ (Definition~\ref{def:groundlens}); not the visible-predicate family of any
map on $\omega$ (Proposition~\ref{prop:nolens}). $\mathcal B_M[x]$ is the Boolean algebra generated
by $\mathcal B_M$ and $x$.
\item[Restriction filter.] For a total $x$, $G_x=\{p:p\subseteq x\}$.
\item[$D_g$, $E_A$, $D^\ast$.] Disagreement with $g$; splitting of $A$; disagreement within some pair
$\{2n,2n+1\}$ (\S\S\ref{ssec:finitebridge}, \ref{ssec:converse}).
\item[Coded closed set.] The branch set $[T]$ of a tree of conditions $T\in M$, computed outside
$M$ (\S\ref{ssec:converse}).
\item[Strict extension.] An enlargement of the algebra of expressible predicates on the same
domain, $\mathcal B_M\subsetneq\mathcal B_M[x]$ (Theorem~\ref{thm:strict}); in the framework's
coarser formulation, a failure of factorization (Theorem~\ref{thm:nonfactor}). It is not an
inclusion of first-order theories.
\item[Package change.] New content obtained by adjoining something the fixed package cannot produce
or express (Corollaries~\ref{cor:packagechange} and \ref{cor:strictcohen}).
\end{description}

\subsection{Hypotheses and regimes}\label{subsec:gate-predicates-list}

\begin{description}[style=multiline,leftmargin=!,labelwidth=0.27\textwidth,labelsep=0.8em,itemsep=3pt]
\item[\textbf{A\_FIN}.] The framework's finite-carrier regime. Here it applies to each
presentation and each element of $\HF$ separately: finite gates, finite transport, least-index
sections and finite diagnostics. It does not license any infinitary commitment.
\item[\textbf{(S)}.] The structural premise: a transitive admitted universe $U$ closed under empty
set, pairs, unions and relatively definable subsets (\S\ref{ssec:box}).
\item[\textbf{C-INF}.] An inductive set in $U$; priced by Theorem~\ref{thm:nonattain}.
\item[\textbf{C-POW}.] For each $x\in U$, an element of $U$ whose members are exactly the subsets of
$x$ lying in $U$; its finite form is Corollary~\ref{cor:finpow}.
\item[\textbf{C-REPL}.] Images in $U$ of sets in $U$ under maps definable in $U$; its finite form
is the transport clause.
\item[\textbf{C-SEL}.] Sections in $U$ for all surjections in $U$; equivalent to Choice in $U$ when
$U$ satisfies $\mathsf{ZF}$ (Corollary~\ref{cor:choice}).
\item[\textbf{J1}.] A countable transitive model $M$ of $\mathsf{ZFC}$, countable externally.
\item[\textbf{J2}.] The classical generic-extension machinery for Cohen forcing, used only to
identify $M[G]$.
\item[Diagnostic gates.] Pass/fail criteria in the computational exhibits. They test encodings and
numerical predictions on bounded cases and do not raise any statement to theorem grade.
\end{description}
